
\subsubsection{6.1 Interpreting the Layer-Type Stratification}

The Conv2d/Linear variance ratio difference (0.637 vs. 0.133) is consistent with three non-exclusive structural explanations.

\textbf{Structural explanation.} Conv2d weight tensors W in R^{c_out x c_in x H x W} have spatial structure (H x W kernel dimensions) that constrains the active GL directions. Linear/FC matrices W in R^{c_out x c_in} are fully dense: the full GL(c_in) x GL(c_out) group acts without structural constraint, providing more variance directions than the discrete permutation group S_{c_in} x S_{c_out}.

\textbf{Training dynamics explanation.} The trajectory finding (ratio approximately 0.49 to 0.28 over 50 epochs) is consistent with gradient descent exploiting GL-type flat-direction reparameterizations during optimization. At initialization, weights drawn from symmetric distributions may have higher permutation orbit density; training progressively moves the weights into non-permutation-orbit directions. This is consistent with the loss landscape literature on flat directions along symmetry orbits.

\textbf{Scale explanation.} Linear layers with large input dimensions have higher-dimensional GL(d) groups, providing more variance directions than S_d (permutation), amplifying GL dominance relative to permutation.

These explanations are not mutually exclusive and are not tested directly in this work.

\subsubsection{6.2 Relation to Prior Literature}

The H-M2 result provides independent empirical support for the theoretical motivation behind GL-invariant trace features in Transformer-NFN [Tran-Viet et al., 2024]. That work proposed tr(WW^T) and tr(W^Q W^{K,T}) precisely because GL symmetry is expected to dominate linear/attention weights. The present measurement of ratio = 0.133 for Linear (FC) layers in CNN Zoo is consistent with this expectation from a different methodological angle. Note, however, that the present measurement covers only CNN Zoo Linear layers; Transformer Zoo attention layers were not directly measured (see Limitation L3).

The Conv2d ratio = 0.637 is consistent with NFN's performance on CNN Zoo (tau > 0.93 using permutation-only equivariance): permutation equivariance captures the dominant variation in convolutional weights in this dataset.

\subsubsection{6.3 Implications for Cross-Architecture Weight Representation Design}

The H-M2 result suggests that a single permutation-based positional encoding is insufficient for cross-architecture WSL because the dominant symmetry differs by layer type. The pre-specified hybrid encoding — permutation orbit-PE for Conv2d layers combined with GL-invariant trace features for Linear and attention layers — is a direct response to this finding and is not post-hoc. The infrastructure developed in this work (orbit projectors, variance decomposer, zoo-scale evaluation pipeline) is directly reusable for implementing and testing the hybrid approach.

Whether GL dominance in CNN Zoo Linear layers extends to Transformer Zoo attention layers is an open measurement question. The inference that it does is plausible but unconfirmed by direct measurement (Limitation L3).

\subsubsection{6.4 Limitations}

\textbf{L1 — Primary performance claim not tested.} The original hypothesis predicted tau_retention >= 0.70 with SANE+orbit-PE. This prediction was not tested. H-M3, which would have tested this claim, was blocked by H-M2's gate failure. The claim is neither confirmed nor refuted; it is empirically unaddressed.

\textbf{L2 — SVHN cross-dataset stability not verified.} SVHN Zoo data was unavailable during H-M2 execution. The pre-specified cross-dataset stability check (|CIFAR ratio - SVHN ratio| < 0.10) was not performed. The CIFAR-10-GS result (n = 1,000 x 50) is based on a single dataset.

\textbf{L3 — Transformer Zoo variance not measured.} The Var_perm/Var_GL ratio for Transformer Zoo checkpoints was not computed. The discussion in Section 6.3 treats the Linear ratio from CNN Zoo as informative about attention layers by inference, but this is not direct measurement.

\textbf{L4 — SVD fallback in H-M2.} A path resolution issue prevented H-M2's orbit_projector.py from importing H-M1's OrbitPEComputer directly. An SVD-based orbit basis was used as a fallback. The two approaches are methodologically equivalent for variance decomposition purposes but represent a minor deviation from the planned implementation design.

\textbf{L5 — No comparison to published tau baselines.} Formal comparison against published NFN and SANE tau values was not conducted. All tau comparisons in the text are from the cited papers.

\subsubsection{6.5 Broader Implications}

The finding that permutation-orbit variance fraction is layer-type-dependent is a general property of weight space geometry that applies to any cross-architecture WSL method. The training trajectory finding (ratio decreasing during training) has a practical implication: early-epoch checkpoints may provide stronger permutation orbit signal than well-trained models. This has potential relevance for model zoo design when permutation-equivariant representations are intended.

